% ID: FIS-OTT-LEN-001
% Titolo: Immagine prodotta da una lente convergente
% Disciplina: Fisica
% Area: Ottica
% Argomento: Lenti
% Sottoargomento: Equazione delle lenti sottili e ingrandimento
% Classe: Seconda liceo scientifico
% Difficolta: 3
% Tipo: problema_numerico
% Risultato: q = 20 cm; h' circa -2,7 cm; immagine reale e capovolta
% Tempo_stimato: 8 min
% Competenze: uso dell'equazione delle lenti, calcolo dell'ingrandimento, interpretazione dell'immagine
% Prerequisiti: lenti convergenti, distanza focale, ingrandimento
% Tag: fisica, ottica, lenti, lente_convergente, ingrandimento
% Fonte: esempio_originale
% Autore: OpenAI Codex
% Licenza: CC-BY-SA-4.0

\begin{esercizio}
Un oggetto alto \(4{,}0\,\mathrm{cm}\) si trova a \(30\,\mathrm{cm}\) da
una lente convergente di distanza focale \(12\,\mathrm{cm}\).

Determina la posizione dell'immagine, la sua altezza e stabilisci se
l'immagine e diritta o capovolta.
\end{esercizio}

\begin{soluzione}
Usiamo l'equazione delle lenti sottili:
\[
\frac{1}{f}=\frac{1}{p}+\frac{1}{q},
\]
dove \(p\) e la distanza dell'oggetto e \(q\) quella dell'immagine.

Ricaviamo \(q\):
\[
\frac{1}{q}=\frac{1}{f}-\frac{1}{p}
=\frac{1}{12}-\frac{1}{30}.
\]

Portando a denominatore comune:
\[
\frac{1}{q}=\frac{5}{60}-\frac{2}{60}
=\frac{3}{60}
=\frac{1}{20}.
\]

Quindi:
\[
q=20\,\mathrm{cm}.
\]

L'ingrandimento vale:
\[
G=-\frac{q}{p}=-\frac{20}{30}=-\frac{2}{3}.
\]

L'altezza dell'immagine e:
\[
h'=Gh=-\frac{2}{3}\cdot 4{,}0\,\mathrm{cm}
\simeq -2{,}7\,\mathrm{cm}.
\]

Il segno negativo indica che l'immagine e capovolta. Poiche \(q\) e positivo,
l'immagine e reale e si forma dalla parte opposta della lente rispetto
all'oggetto.
\end{soluzione}
